% packages.tex

\usepackage[utf8]{inputenc}   % Permette di leggere i caratteri accentati e speciali da tastiera
\usepackage[T1]{fontenc}  % Usa l'encoding T1
\usepackage[english]{babel}
\usepackage{listings}
\usepackage{xcolor}  % Per usare i colori
\usepackage{lipsum}
\usepackage{scrhack} % resolve KOMA-Script warnings


% -- MATH -- %

\usepackage{amsmath}
\usepackage{cancel}
\usepackage{amsthm} % math theorem, proofs, ...
\usepackage{thmtools}  % custom theorem styles

\declaretheoremstyle[
    headfont=\bfseries,       
    bodyfont=\normalfont,                  
]{base_th_style}

\declaretheoremstyle[
    style = base_th_style,    
    numbered=no               
]{thmstyle}

\declaretheoremstyle[
    style=base_th_style,    
    numbered=yes               
]{lemstyle}

\declaretheoremstyle[
    style=base_th_style,    
    numbered=yes               
]{defstyle}

\declaretheorem[style=thmstyle, name=Theorem]{theorem*}
\declaretheorem[style=thmstyle, name=Property]{property*}
\declaretheorem[style=lemstyle, name=Lemma]{lemma}
\declaretheorem[style=defstyle, name=Definition]{definition}
\renewenvironment{proof}{{\noindent\bfseries Proof.}}{\hfill$\square$\medskip}

\newcommand{\R}{\mathbb{R}}
\newcommand{\N}{\mathbb{N}}
\newcommand{\Z}{\mathbb{Z}}
\newcommand{\Q}{\mathbb{Q}}
\newcommand{\C}{\mathbb{C}}
\newcommand{\ve}{\varepsilon}
% --  -- %

\usepackage{import}
\usepackage{caption}

\usepackage{graphicx}
\usepackage{float}
\newcommand{\extlink}[2]{%
    \href{#1}{\color{lightblue}\uline{#2}}%
}
\usepackage{soul} % for cut text
\usepackage{csquotes}
\usepackage{fancyvrb} % for verbatim environment
\usepackage[backend=biber, style=numeric, sorting=none]{biblatex}
\addbibresource{appendices/ref.bib}

\usepackage[backend=biber, style=numeric, sorting=none]{biblatex}
\addbibresource{appendices/ref.bib}

\usepackage{bm} % bold math block

\usepackage{lmodern}          % Font di alta qualità che supporta meglio gli accenti rispetto al default
\usepackage{textcomp}         % Supporta simboli aggiuntivi

\usepackage{framed}
\usepackage{enumitem}
\usepackage{multirow}
\usepackage{multicol}
\usepackage{array}
\usepackage{tabularx} % for better tables
\usepackage[table]{xcolor} % for coloring table rows
\usepackage{makecell}

% table in multiple pages
\usepackage{ltablex}
\usepackage{booktabs}
\keepXColumns

% particular schemes
\usepackage{tikz}
\usepackage{pgfplots}
\pgfplotsset{compat=1.18}
\usetikzlibrary{shapes, arrows.meta, positioning}
\usetikzlibrary{calc} % for coordinate calculation

% -------------------------------- FSM --------------------------------%
\usetikzlibrary{decorations.pathreplacing} % for braces
\usetikzlibrary{automata, positioning, arrows} % for FSM diagrams
\usetikzlibrary{backgrounds} % for background layers

% States and transitions style
\tikzset{
    fsm_style/.style={
        >=stealth, 
        node distance=4cm, 
        on grid, 
        auto,
        thick,
        every state/.style={minimum size=1.5cm}
    }
}
% -------------------------------- FSM --------------------------------%

\usepackage[most]{tcolorbox}
\usepackage{fontawesome5}
\usepackage{xcolor}
\usepackage{url}
\usepackage{hhline}
